\PassOptionsToPackage{unicode}{hyperref} 
\documentclass{beamer}

% Налаштування мови та шрифтів
\usepackage[T2A]{fontenc}
\usepackage[utf8]{inputenc}
\usepackage[ukrainian]{babel}
\usepackage{graphicx}
\usefonttheme[onlymath]{serif}

% =========================
% 🔧 ВИБІР ТЕМИ (ЗМІНИ ТУТ)
% =========================

\newcommand{\mytheme}{Berlin} 
% Спробуйте замінити на: Berlin, Madrid, Warsaw, Frankfurt, AnnArbor

\usetheme{\mytheme}

% 🔧 Кольорова тема
\usecolortheme{seahorse} 
% Спробуйте: seahorse, whale, beetle, orchid

% =========================

\title{Тема: \mytheme}
\author{Миколаєнко Дмитро}
\date{\today}

\begin{document}

% Титульний слайд
\begin{frame}
    \titlepage
\end{frame}

% Слайд 1
\section{Приклад}
\begin{frame}{Список у темі \mytheme}
    \begin{itemize}
        \item<1-> Перший пункт
        \item<2-> Другий пункт
        \item<3-> Третій пункт
    \end{itemize}
\end{frame}

% Слайд 2
\begin{frame}{Математика}
    Формула площі квадрата:
    \[
    S = a^2
    \]
\end{frame}

% Слайд 3
\begin{frame}{Графіка}
    \begin{center}
        % Пакунок graphicx дозволяє вставити стандартну заглушку
        \includegraphics[width=0.4\textwidth]{example-image}
    \end{center}
\end{frame}

% Слайд 4
\begin{frame}{Висновок}
    \begin{itemize}
        \item[\checkmark] **Найкраща тема:** [Berlin]
        \item[\checkmark] **Найкращий колір:** [seahorse]
        \item[\checkmark] **Моє спостереження:** [Обрана тема і колір є світлими, що забезпечує зручність сприйняття інформації та покращує читабельність тексту.]
       
    \end{itemize}
\end{frame}

\end{document}
